\documentclass{ximera}

\input{../../preamble.tex}

\title{Categories}

\begin{document}

\begin{abstract}
templates
\end{abstract}
\maketitle





We are building a library of the elementary functions.  The idea is to use the library to list characteristics, features, and aspects of all functions within each category.  

That way, if we can identify the type of function we have, then we get free information when analyzing functions. 

The category becomes our reasoning. 



\begin{center}

\textbf{\textcolor{red!70!black}{These are ``CAN'' questions.}} 

\end{center}




\textbf{\textcolor{purple!85!blue}{CAN}} the formula we are given be rewritten as one of the official standard forms for each category? 






\section*{Official Templates of Elementary Functions}


These elementary function categories are our \textbf{first choice}.  If a function can be represented by one of these standard forms, then we want to describe the function as one of these elementary functions.  That gives us the most information. 

\begin{formula} \textbf{\textcolor{blue!55!black}{Power Functions}} 

A power function is any function that \textbf{\textcolor{purple!85!blue}{CAN}} be represented with a formula of the form

\[   pow(x) = k \, x^r      \]

where $k$ and $r$ are real numbers.




\end{formula}












Polynomial functions are sums of power functions with powers that are nonnegative integers (whole numbers).


\begin{formula} \textbf{\textcolor{blue!55!black}{Polynomial Functions}} 

A polynomial function is any function that \textbf{\textcolor{purple!85!blue}{CAN}} be represented with a formula of the form

\[   poly(x) =  a_n x^n + a_{n-1} x^{n-1} + \cdots + a_3 x^3 + a_2 x^2 + a_1 x^1 + a_0 x^0      \]

where the $a_k$ are real numbers and $a_n \ne 0$. 



\textbf{Special Polynomials}
\begin{itemize}
\item \textbf{\textcolor{blue!55!black}{Constant:}} $C(x) = C$ 
\item \textbf{\textcolor{blue!55!black}{Linear:}} $L(x) = A \, x + B$ 
\item \textbf{\textcolor{blue!55!black}{Quadratic:}} $Q(x) = A \, x^2 + B \, x + C$ where $A \ne 0$
\end{itemize}

\end{formula}











\begin{formula} \textbf{\textcolor{blue!55!black}{Rational Functions}} 

A rational function is any function that \textbf{\textcolor{purple!85!blue}{CAN}} be represented with a formula of the form

\[ rat(x) =   \frac{ a_n x^n + a_{n-1} x^{n-1} + \cdots + a_3 x^3 + a_2 x^2 + a_1 x + a_0  } { b_m x^m + b_{m-1} x^{m-1} + \cdots + b_3 x^3 + b_2 x^2 + b_1 x + b_0 }   \]



where the $a_k$ and $b_k$ are real numbers and $a_n \ne 0$ and $b_m \ne 0$.





\end{formula}

















\begin{formula} \textbf{\textcolor{blue!55!black}{Radical/Root Functions}} 

A radical or root function is any function that \textbf{\textcolor{purple!85!blue}{CAN}} be represented with a formula of the form  

\[   rad(x) = A \sqrt[n]{B x + C} + D =  A (B \, x + C)^{\tfrac{1}{n}} + D    \]

where the $A$, $B$, $C$, and $D$ are real numbers and $A \ne 0$ and $B \ne 0$.

\end{formula}














\section*{Operations}


(1) The elementary function categories above are our \textbf{first choice}.  If a function can be represented by one of these standard forms, then we want to describe the function as one of these elementary functions.  That gives us the most information. 


(2) If a function cannot be identified as one of these elementary functions, then our \textbf{second choice} is an operation. 





\begin{center}

\textbf{\textcolor{red!70!black}{These are ``IS'' questions.}} 

\end{center}




\textbf{\textcolor{purple!85!blue}{IS}} the formula we are given written as one of the official standard forms for an operation? 

NOT can it be written, but is it written as an operation?





Our operational interpretation of every expression is through the \textbf{\textcolor{purple!85!blue}{Order of Operations}}. 

Every mathematical expression involving mathematical operations is either

\begin{itemize}
	\item a \textbf{Constant Multiple} : A f(x);
	\item a \textbf{Sum} : f(x) + g(x);
	\item a \textbf{Difference} : f(x) - g(x);
	\item a \textbf{Product} : f(x) g(x); or 
	\item a \textbf{Quotient} : \frac{f(x)}{g(x)};
\end{itemize}

\textbf{Note:} A constant multiple is techinally a product, with the first factor a constant.

Each mathematical expression is one of these and only one of these.  The Order of Operations tells us which.













\section*{Composition}


(1) These elementary function categories above are our \textbf{first choice}.  If a function can be represented by one of these standard forms, then we want to describe the function as one of these elementary functions.  That gives us the most information. 


(2) If a function cannot be identified as one of these elementary functions, then our \textbf{second choice} is an operation. 


(3) If a function is not an elementary function and it is not one of our four operations, then our \textbf{third choice} is a \textbf{\textcolor{purple!85!blue}{composition}}.  We study composition in great depth in Calculus. 

(0) There is a fourth choice and that is a piecewise defined function.  Usually, these are easy to identify, because they have a big curly brace in front of several formulas. 

(4) If a function is not any of these, then it is weird.

Incidentally, we like weird functions.





































\begin{center}
\textbf{\textcolor{green!50!black}{ooooo-=-=-=-ooOoo-=-=-=-ooooo}} \\

more examples can be found by following this link\\ \link[More Examples of Elementary Functions]{https://ximera.osu.edu/csccmathematics/precalculus1/precalculus1/elementaryLibrary1/examples/exampleList}

\end{center}





\end{document}
